\chapter{Formulas for variance and concentration}
In this chapter, we introduce the  critical exponents of the model and relations between them. The starting points are precise formulas for  fluctuations  of the partition function (variance and large deviations).
%%
\section{The critical exponents}
 There are different ways of defining the critical exponents, see for example \cite{Chatterjee-scaling,LiceaNewPiza}. We will not enter the finest details, and we stay at an intuitive level.
 Although it is not clear that these definitions are all equivalent, the main idea is that the critical exponents are two reals $\xi^{\perp}$ and $\xi^{\parallel}$ such that
\begin{equation}
\sup_{0\leq s \leq t} |B_s| \approx t^{\xi^{\perp}(d)} \quad \text{and} \quad \ln Z_t - \IP[\ln Z_t] \approx t^{\xi^{\parallel}(d)} \quad \text{as} \  t\to\infty.
\end{equation}
The "wandering exponent" $\xi^{\perp}$ is the exponent for the asympotic transversal (or "perpendicular") fluctuations of the path, with respect to the time axis.
%Under the P2P measure, it measures the deviation of the polymer path from the 
 The polymer is said to be \emph{diffusive} when $\xi^{\perp} = 1/2$ (as for the brownian motion), and it is said to be \emph{super-diffusive} when $\xi^{\perp} > 1/2$. One of the conjectures in polymers is that diffusivity should occure in weak disorder, while super-diffusivity should take place in the strong disorder setting. The number $\xi^{\parallel}$ denotes the critical exponent for the longitudinal fluctuation of the free energy.

The study of these exponents goes beyond the polymer framework. 
The reason is that they are expected to take the same value in many different  statistical physics models describing growth phenomena.  In dimension $d=1$ 
this family is called the KPZ universality class (see Section \ref{sec:IR}). It is conjectured in the physics literature \cite{KrugSpohn} that the two exponents should depend on one other, in the way that
\begin{equation}
\xi^{\parallel}(d) = 2 \xi^{\perp}(d) - 1, \forall d\geq 1
\end{equation}
Under a certain definition of the exponents, the relation was proved by Chaterjee \cite{Chatterjee-scaling} for first-passage percolation. Auffinger and Damron were able to simplify Chatterjee's proof and extend the result to directed polymers \cite{AuffingerDamron-sc,AuffingerDamron13}.

In dimension $d=1$, it is conjectured that $\xi^{\perp}=2/3$ and $\xi^{\parallel}=1/3$ for any positive $\b$. For now, this has only been proven for solvable models of polymers: Sepp\"al\"ainen's discrete log-gamma polymer \cite{Sepp12}, O'Connell-Yor semi-discrete polymer \cite{MorenoSeppValko,SeppValko10}, and also for the the KPZ polymer \cite{BalaszQuastelSeppalainen11}. 


In dimension $d \geq 2$, essentially nothing is known. Let's simply mention the (rough) bounds
$$
0 \leq \xi^\parallel \leq 1/2\,,\quad 1/2 \leq \xi^\perp \leq 3/4\,,
$$
where the last one will be proved in section \ref{ch:CM}.

A way to approach $\xi^{\parallel}$ is to consider the variance of  $\ln Z_t$.
In what follows, we give a formula to express the variance of $\ln Z_t$ in terms of a stochastic integral, which is obtained through a Clark-Ocone type martingale representation.

\section{The Clark-Ocone representation}
It is a consequence of It\^o's work on iterated stochastic integrals \cite{ito1951multiple} any that square-integrable functionals of the Brownian motion can be written as the sum of a constant and an It\^o integral. In \cite{clark1970representation}, Clark extended this result to a wider range of functionals, and showed that any martingale that is measurable with respect to the Brownian motion filtration, could be represented as a stochastic integral martingale. Clark was also able to compute the integrand of the representation, for a special class of functionals. Ocone then showed \cite{ocone1984malliavin} that this computation was linked to Malliavin's calculus, and generalised this idea to a larger class of functionals.

Such representations - called Clark-Ocone representations - also exist in the framework of functionals of a Poisson processes. 
%Starting in the end of the 90's, several papers
%\cite{picard1996formules,wu2000new,ane2000logarithmic}
% showed that this was the case for particular choices of Poisson processes. In \cite{LAST20111588}, Last and Penrose proved a Clark-Ocone type martingale representation formula for general Poisson processes. In their paper, they also treated the particular case of a Poisson process $\om$ on $\mathbb{R}_+\times \mathbb{R}^d$. 
Denote by $\om_{s-}$ the restriction of $\om$ on $[0,s)\times\mathbb{R}^d$, and consider the derivative operator
\begin{equation}
D_{(s,x)}F(\om) := F(\om + \delta_{s,x}) - F(\om).
\end{equation}
We have:
\begin{theorem} \textup{\cite[theorem 3.1]{LAST20111588}}
Let $F=F(\om)$ be a functional of the Poisson process, such that $\IP[F^2]<\infty$. Then,
\begin{equation}
\IP \int \IP[D_{(s,x)}F(\om)|\om_{s-}]^2 \rmd s \rmd x < \infty,
\end{equation}
and we have for all $u\geq 0$, that $\IP$-a.s.
\begin{equation}
\IP[F(\om)|\om_u] = \IP[F(\om)] + \int_{[0,u]\times \mathbb{R}^d} \IP[D_{(s,x)} F(\om) | \om_{s-}] \bar \om (\rmd s \rmd x).
\end{equation}
\end{theorem}
This proves that the square integrable martingale $(\IP[F(\om)|\om_u])_{u\geq 0}$ admits a stochastic integral martingale representation, with predictable integrand $\IP[D_{(s,x)} F(\om) | \om_{s-}]$.

\section{The variance formula}
To lighten the writing, we will denote by $\mathcal{G}_{s-}$ the sigma-field generated by $\om_{s-}$ and $\IP^{\mathcal{G}_{s-}}$ will stand for the expectation knowing $\mathcal{G}_{s-}$.

Using Jensen's inequality and Tonelli's theorem, it is easy to check that $\ln Z_t$ is a square integrable function of $\om$. Hence, the process $\big(\IP[\ln Z_t|\om_u]\big)_{u\in[0,t]}$ is a martingale which admits a Clark-Ocone type representation:
\begin{equation} \label{eq:ClarKRepresentationLogZt}
\IP[\ln Z_t|\om_u] = \IP\left[\ln Z_t\right] + \int_{[0,u]\times\mathbb{R}^d} \IP^{\mathcal{G}_{s-}}\left[D_{(s,x)}\ln Z_t\right] \bar{\om}(\rmd s \rmd x),
\end{equation}
where
\[D_{(s,x)}F(\om) = \ln \frac{P\left[e^{\b\om(V_t)}e^{\b \chi_{s,x}}\right]}{Z_t}  = \ln \left(1+\lambda \gibbs[\chi_{s,x}]\right).\]
As a consequence, one can express the variance of $\ln Z_t$ via \eqref{eq:ClarKRepresentationLogZt}, using the formula for the variance of a Poisson integral \eqref{eq:MartingaleVarianceFormula}, and find that
\begin{equation} \label{eq:varFormula}
\mathrm{Var} (\ln Z_t) = \IP \int_{[0,t]\times\mathbb{R}^d} \IP^{\mathcal{G}_{s-}}\left[\ln \left(1+\lambda \gibbs[\chi_{s,x}]\right)\right]^2 \nu \hspace{0.5mm} \rmd s \rmd x.
\end{equation}

This variance formula leads us to the following theorem:
\begin{theorem} \label{th:var}
\begin{enumerate}[label=(\roman*)]
  \item The following lower and upper bounds on the variance hold:
  \begin{align}
\mathrm{Var} (\ln Z_t) &\geq c_-^2 \IP \int_{[0,t]\times\mathbb{R}^d} \IP^{\mathcal{G}_{s-}}\big[ \gibbs [\chi_{s,x}] \big]^2 \, \nu \hspace{0.5mm} \rmd s \rmd x,\\
\mathrm{Var} (\ln Z_t) &\leq c_+^2 \IP \int_{[0,t]\times\mathbb{R}^d} \IP^{\mathcal{G}_{s-}}\big[ \gibbs [\chi_{s,x}] \big]^2 \, \nu \hspace{0.5mm} \rmd s \rmd x, \label{eq:IneqVar}
  \end{align}
where $c_- = 1 - e^{-|\b|}$ and $c_+ = e^{|\b|}-1$.

In particular,
  \begin{equation}
  \mathrm{Var} (\ln Z_t) \leq c_+^2 \, t \nu \, \IP[J_t]. \label{eq:var_bound}
  \end{equation} 
  \item Letting $c=\nu c_+^2 \exp(c_+)$, the following concentration estimate holds:
  \begin{equation} \label{eq:concentrationlogZt}
  \IP\left(\big|\ln Z_t - \IP[\ln Z_t]\big| > u\right) \leq 2\exp\left(-\frac{1}{2}(u\wedge\frac{u^2}{ct})\right).
  \end{equation}
\end{enumerate}
\end{theorem}
\begin{remark} Recalling Theorem \ref{th:asymptOfJt}, the inequality \eqref{eq:var_bound} suggests that the variance should be smaller in weak disorder than in strong disorder. It also shows that for all $d\geq 1$, we have $\xi^{\parallel}(d)\leq 1/2$.
\end{remark}
\begin{proof}
The two first bounds on the variance are a consequence of the fact that, for all $u\in [0,1]$, we have 
\begin{equation}\label{eq:LogBound}
c_- u \leq |\ln(1+\lambda u)| \leq c_+ u.\end{equation} Then, apply Jensen's inequality to the conditional expectation in the right-hand side of \eqref{eq:IneqVar} and use Fubini's theorem such that
\begin{align*}
\mathrm{Var} (\ln Z_t) &\leq c_+^2 \IP \int_{[0,t]\times\mathbb{R}^d} \IP^{\mathcal{G}_{s-}}\big[ \gibbs [\chi_{s,x}] \big]^2 \, \nu \hspace{0.5mm} \rmd s \rmd x \\
&\leq c_+^2 \IP \int_{[0,t]\times\mathbb{R}^d} \gibbs [\chi_{s,x}]^2 \, \nu \hspace{0.5mm} \rmd s \rmd x \\
& = c_+^2 t \nu \, \IP[J_t],
\end{align*}
by definition of $J_t$. This completes the proof of (i).
To prove \eqref{eq:concentrationlogZt}, we first denote by $Y_{t,u}$ the mean-zero martingale part appearing in \eqref{eq:ClarKRepresentationLogZt}, i.e.
\[Y_{t,u} := \int_{[0,u]\times\mathbb{R}^d} \IP^{\mathcal{G}_{s-}}\ln \left(1+\lambda \gibbs[\chi_{s,x}]\right) \bar{\om}(\rmd s \rmd x).\]
Then, letting $\varphi(v)=e^v - v - 1$ and $a\in [-1,1]$, we define $(M_{t,u})_{u\in[0,t]}$ as the exponential martingale associated to $(Y_{t,u})_{u\in[0,t]}$:
\[M_{t,u}=\exp\left(aY_{t,u}-\int_{[0,u]\times \mathbb{R}^d}  \varphi\left(a \cdot\IP^{\mathcal{G}_{s-}}\ln \big(1+\lambda \gibbs[\chi_{s,x}]\big)\right) \, \nu \hspace{0.5mm} \rmd s \rmd x\right).\]
By \eqref{eq:LogBound} and the observations that $\chi$ is less than $1$ and that $|\varphi(v)|\leq e^{|v|} v^2/2$ for all $v$, we have for $a\in[-1,1]$
\begin{align*}\left| \int_{[0,t]\times \mathbb{R}^d}  \varphi\left(a \cdot\IP^{\mathcal{G}_{s-}}\ln \left(1+\lambda \gibbs[\chi_{s,x}]\right)\right) \, \nu \, \rmd s \rmd x \right| 
&\leq e^{c_+} \frac{c_+^2 a^2}{2} \int_{[0,t]\times\mathbb{R}^d} \IP^{\mathcal{G}_{s-}}\big[ \gibbs [\chi_{s,x}] \big]^2 \, \nu \hspace{0.5mm} \rmd s \rmd x\\
&\leq c \frac{a^2}{2} \int_{[0,t]\times\mathbb{R}^d} \IP^{\mathcal{G}_{s-}}\big[ \gibbs [\chi_{s,x}] \big] \, \rmd s \rmd x\\
&= c \frac{a^2}{2}t,
\end{align*}
where $c = \nu c_+^2 e^{c_+}$ and where the second inequality was obtained using Jensen's inequality. 

If one denotes by $b_{t,u}$ the integral term in the definition of $M_{t,u}$, we just showed that $b_{t,t} \leq c a^2t/2$, so by Markov's inequality and the martingale property, we obtain that
\begin{equation}
\IP\big(\ln Z_t - \IP[Z_t] > u\big) = \IP\big(M_{t,t} > \exp(au-b_{t,t})\big)
\leq \exp(c a^2t/2 - au).
\end{equation}
This implies \eqref{eq:concentrationlogZt} after minimizing the bound for $a\in[-1,1]$, and repeating the same procedure for the lower deviation.
\end{proof}

From the concentration estimate, one can derive the following almost sure behavior:
\begin{corollary} \label{cor:var1}
For all $\varepsilon>0$ and as $t\to \infty$,
\begin{equation} \label{eq:OrderOfCenteredlogZt}
\ln Z_t - \IP[\ln Z_t] = \mathcal{O}\left(t^{\frac{1+\varepsilon}{2}}\right), \quad \IP\text{-a.s.} 
\end{equation}
\end{corollary}
\begin{proof}
Equation \eqref{eq:concentrationlogZt} implies that for large enough $k\in\mathbb{N}$,
\begin{equation}
\IP\left[|\ln Z_k - \IP \ln Z_k|> k^{\frac{1+\varepsilon}{2}}\right] \leq 2\exp\big(-\frac{t^{\varepsilon}}{2c}\big),
\end{equation}
which is summable. By Borel-Cantelli lemma, we obtain that,  $\IP$-almost surely,
$$|\ln Z_k - \IP[\ln Z_k]| \leq k^{\frac{1+\varepsilon}{2}}\qquad {\rm for\ } k {\rm \ large\ enough.}
$$ 
To extend this to any $t\geq 0$ and prove \eqref{eq:OrderOfCenteredlogZt}, it suffices to apply the next lemma.
\end{proof}
\begin{lemma} Let $h>0$. For all $0\leq s \leq h$,
\begin{equation} \label{eq:BoundOnZt+s-Zt}
-c_+ \delta_t(h) \leq \ln Z_{t+s} - \ln Z_t \leq c_+\delta_t(h),
\end{equation}
where
\[\delta_t(h)=\int_{[t,t+h]\times\mathbb{R}^d} P_{s-}^{\b,\om}[\chi_{s,x}] \om(\rmd s \rmd x),\]
is such that for any $\varepsilon > 0$,
\begin{equation}\label{eq:deltah}\delta_t(h)=\mathcal{O}\big(t^{\frac{1+\varepsilon}{2}}\big), \quad \IP\text{-a.s.}\end{equation}
\end{lemma}
\begin{proof} We get from the integral writing of $\ln Z_t$ \eqref{eq:logZtIntRepresentation} that
\[\ln Z_{t+s} - \ln Z_t = \int_{[t,t+h]\times\mathbb{R}^d} \om(\rmd s \rmd x) \ln \big(1+\lambda P_{s-}^{\b,\om} [\chi_{s,x}]\big).\]
Hence, \eqref{eq:BoundOnZt+s-Zt} is simply obtained with \eqref{eq:LogBound}.

Now, introduce the martingale
\[M_t = \int_0^t \om(\rmd s \rmd x) P_{s-}^{\b,\om}[\chi_{s,x}] - \nu r^dt,\]
which has bracket ${\langle M \rangle_t}=\nu\int_0^t I_s \rmd s \leq \nu r^d t$. Note that $\delta_t(h) = M_{t+h}-M_t + h$, so that \eqref{eq:deltah} is thus a consequence of the martingale properties
\eqref{eq:BracketImpliesCV} and
\eqref{eq:InfiniteBracket}, since in the case where $\langle M \rangle_\infty$ is infinite, then $|M_t| = o\big(\langle M \rangle_t^{\frac{1+\varepsilon}{2}}\big)$ as $t\to\infty$.
\end{proof}
%%%
%[??? corollaire a garder? remarque a supprimer]
In order to illustrate the general strategy, we now mention a consequence of theorem \ref{th:var} (i). 


\begin{corollary} \label{cor:var2} Let $\b\neq 0$, $\xi$ and $C>0$. There exists a constant $c_1 = c_1(d,C)\in(0,\infty)$, such that 
\begin{align*} \liminf_{t\to\infty}t^{-(1-d\xi)}\mathrm{Var}(\ln Z_t) &\geq c_1 \liminf_{t\to\infty} \inf_{0\leq s \leq t} \left(\IP \gibbs\left(|B_s|\leq C + Ct^\xi\right)\right)^2\\
&\geq c_1 \liminf_{t\to\infty} \left(\IP \gibbs\left(\sup_{0\leq s \leq t} |B_s| \leq C + Ct^\xi\right) \right)^2.
\end{align*}
\end{corollary}
The result suggests that 
$$
\chi(d) \geq \frac{1-d\xi(d)}{2}.
$$
For more details on the result and the proof, we refer the reader to Corollary 2.4.3 in \cite{CY05}.
%\begin{remark}
%This suggests in particular that for $\b\neq 0$,
%\begin{equation}\nn
%\chi(d) \geq \frac{1-d\xi(d)}{2}.
%\end{equation}
%Note that it is not the case when $\b = 0$ [??? since $\chi(d)=-\infty$] , and that the expected relation  would be [??? ref relation entre exposants, lien entre les deux?]
%\end{remark}
%\begin{proof}
%
%\end{proof}
